\scanpage{3}
\mspage{1}

$X$ espace topologique $\Longrightarrow$ \inserted{pré}ordonné par \uline{spécialisation} : $\unsure{\operatorname{Ord}(X)}$
\[
x\leq y\quad\overset{\mathrm{déf}}{\Longleftrightarrow}\quad\overline{\{x\}}\subset\overline{\{y\}}
\]
(ordonné ssi $X$ satisfait l'axiome de séparation $(T_0)$)

$I$ ens \inserted{pré}ordonné $\Longrightarrow$ espace top., fermés les $A\subset I$ t.q. $x\in A,y\leq x\Rightarrow y\in A$
ouverts les ens. fermés dans $I^\circ$ (opposé à $I$)
base d'ouverts : les ens $I_{\geq x}=\widetilde{\{x\}}$ ($\sim$ désigne l'adh. pour $I^\circ$)
\[
\operatorname{Esp}(I)
\]
\margin{\deleted{\unsure{lien foncteurs}}}

\uline{NB1} : Si $I,J$ préordonnés, $f:I\longrightarrow J$
une application, alors $f$ croissante ssi
$f$ continue.
\[
\overline{\{x\}}=I_{\leq x}=\{y\in I\mid y\leq x\}\qquad
\begin{aligned}
\text{donc }x\leq y&\Longleftrightarrow\overline{\{x\}}\subset\overline{\{y\}}\\
&\Longleftrightarrow x\in\overline{\{y\}}
\end{aligned}
\]
\inserted{Plus gén. $I\longrightarrow X$ continue ssi elle est croissante}
\margin{\deleted{\unsure{lien}}}

\uline{NB. 2} Une topologie sur un ens $I$ est associée à
un préordre (\unsure{néc.} unique, \unsure{ce sera celui de}
la spécialisation) ssi les réunions quelc. de
fermés sont fermés $\Longleftrightarrow$ intersections d'ouverts
sont ouverts (i.e. $\exists$ une ``topologie opposée'' !)

\unsure{Ceci} : deux foncteurs
\diagram{\begin{tikzcd}[column sep=large]
\mathcal E=(\operatorname{Esp}) \arrow[r,shift left,"\operatorname{Ord}"] \arrow[d,phantom,"\in"{rotate=90}] & (\operatorname{Préord})=\mathcal O \arrow[l,shift left,"\operatorname{Esp}"] \arrow[d,phantom,"\in"{rotate=90}]\\
X & I
\end{tikzcd}}
avec
\[
\operatorname{Ord}\circ\operatorname{Esp}=\operatorname{id}_{\unsure{\mathcal E}}
\qquad
\operatorname{Esp}\circ\operatorname{Ord}\longrightarrow\operatorname{id}_{\unsure{\mathcal O}}
\]
car l'application $\operatorname{id}_X:X\longrightarrow X$
est continue de $\operatorname{Esp}\operatorname{Ord}X$
dans $X$ (et est $\simeq$ ssi $X$
satisfait la cond. du NB2
\[
\operatorname{Hom}_{\mathcal E}(\operatorname{Esp}(I),X)\xrightarrow{\sim}
\operatorname{Hom}_{\mathcal O}(I,\operatorname{Ord}(X))\qquad\text{adjonction}
\]
$I\longrightarrow X$ est continue ssi elle est croissante

Considérons maintenant $\varphi:X\longrightarrow I$, i.e. une pseudo-partition : $X=\bigcup_{i\in I}X_i^*$
(les $X_i$ mutuellement disjoints). \inserted{Posons $X_J^*=\varphi^{-1}(J)$.} \deleted{continue ssi}
$\varphi^{-1}(I_{\geq i})$ ouverts
 i.e. ssi les $X_{\widetilde i}^*=\bigcup_{j\geq i}X_j^*$ sont ouverts, i.e. ssi les
$X_{\complement\widetilde i}^*=\bigcup_{\substack{j\in I\\j\not\geq i}}X_j^*$
sont fermés. Il faut pour ceci que les
$X_i\overset{\mathrm{déf}}=X_{\overline i}^*=\bigcup_{j\leq i}X_j$
soient fermés, et cette condition est suffisante
si $I$ est fini (car tt fermé est alors $\bigcup$ finie de fermés de la forme $\overline i$)
plus généralement si $\varphi$ est ``loct. finie'' i.e. loc. à ens. de val. fini
(\uline{NB} Si $I'\subset I$, la top. de $I'$ associée au préordre induit est la top. induite
par $I$)

\margin{\deleted{\illegible{3} $\unsure{x\leq y}\Longleftrightarrow\overline{\{x\}}\subset\overline{\{y\}}$ \unsure{en}}}
\margin{* \unsure{comme on veut que la top.}
\unsure{quotient en $I$ implique} topo, la
question est mal posée -- plutôt :
quelles conditions supplémentaires sur
$\varphi$ i.e. les $X_i^*$ imposer pour
avoir unicité ?}

Dans quels cas l'\inserted{pré}\uline{ordre} de l'ens $I$ est-il déterminé
par l'application $\varphi$ i.e. la famille $(X_i)_{i\in I}$ (dans le
cas $\varphi$ continue) ? * \inserted{plus précisément par $(X_i^*)_{i\in I}$ donné}
\deleted{\illegible{3}}
quand \deleted{\unsure{tout}} existe-t-il un ordre sur $I$ qui rende $\varphi$
continue, et quand cet ordre est-il unique ?
Pour l'unicité, il est clair que si $X=\varnothing$ il
faut que $\operatorname{card}I\leq1$, et si $X\ne\varnothing$ que $\varphi$ soit
surjective i.e. les $X_i^*\ne\varnothing$. Supposons donc $\varphi$ surj. i.e.
$X_i^*\ne\varnothing$. On aimerait que la top. de $I$ soit la \uline{top.}
\uline{quotient} i.e. que $X_J$ fermé ssi $J$ fermé en $I$.
Cette topologie sur $I$ est bien telle que les réunions
de fermés est fermé \inserted{(cas $\varphi$ loc. finie)} donc est associée à un ens
\unsure{ordonné de} $I$. Pour cet ordre, on aura donc
$i\leq j$ ssi \deleted{$\unsure{x\in}$} $\forall J\subset I$ tel que $X_J^*$ fermé et $j\in J$, on
a $i\in J$. À quelles conditions cette préordre est-il
un ordre (NB ``en général'' ce préordre aura tte
\unsure{chance} d'être chaotique !) Il faut et s. que
les $\{i\}$ soient localement fermés dans $I$, ce
qui implique que les $X_i^*$ sont loct. fermés. \deleted{\unsure{D'ailleurs}}
\deleted{\unsure{réciproquement}} il f. et s. que \inserted{$\forall i\in I$} l'on ait \deleted{$\unsure{\{i\}}$}
$\{i\}=\overline i\cap\widetilde i$ ce qui équivaut à \deleted{$\unsure{\forall i\in I}$}
$X_{\overline i}^*\cap X_{\widetilde i}^*=X_i^*$,
où $X_{\overline i}^*$ est le plus petit fermé de $X$ \inserted{$\supset X_i^*$} ``saturé'' par la
relation d'équivalence loc. finie assoc., et $X_{\widetilde i}^*$ est le
plus petit ouvert saturé $\supset X_i^*$. On pourrait appeler
une ``\uline{stratification fidèle}'' de $X$ une relation d'équiv.
loc\unsure{t.} finie telle que l'ens préordonné quotient
$I$ soit \uline{ordonné} -- ils correspondent à une famille
$(X_i^*)_{i\in I}$, partitions ensembliste de $X$ \inserted{indicée par $I$ ordonné}
\deleted{\illegible{4}} \unsure{tels que}
\unsure{disons et} telle que
Pour $J\subset I$, on ait $X_J$ fermé ssi $J$ fermé.
donc telle que a) les $X_i\overset{\mathrm{déf}}=\bigcup_{j\leq i}X_j^*$ fermés,
\inserted{i.e. $J$ fermé $\Rightarrow X_J$ est \unsure{fermé} ($\Rightarrow$ les $X_i^*$ loc. f.)}
b) Si $X_J$ est fermé, alors $J$ l'est.
